\begin{figure*}[hptb]
    \centering
    \captionsetup{font=footnotesize}
    \begin{minipage}[b]{0.48\textwidth}
        \centering
        \includegraphics[width=\textwidth]{plots_gpplus_logscale_comparison_paper/buckling_xdim4_gpplus_comparison_RRMSE.pdf}
        \subcaption{RRMSE}
        \label{fig logscale buckling rrmse}
    \end{minipage}
    \hfill
    \begin{minipage}[b]{0.48\textwidth}
        \centering
        \includegraphics[width=\textwidth]{plots_gpplus_logscale_comparison_paper/buckling_xdim4_gpplus_comparison_NIS.pdf}
        \subcaption{NIS}
        \label{fig logscale buckling nis}
    \end{minipage}
    \caption{Evaluations based on RRMSE and NIS (Buckling $D_x=4$) with and without log-scale target preprocessing: Noise is either $\varepsilon \sim \mathcal{N}(0, (0.002\times c)^2)$ or $\varepsilon \sim \mathcal{N}(0, (0.08\times c)^2)$. GPs are used on CPU while TabPFN is used on GPU. All the GPs are trained with the \textit{default} settings of GP+ and GPyTorch packages. GP+ (log) and GPyTorch (log) denote models trained after a log transform of the response, with metrics evaluated on the original scale.}
    \label{fig logscale buckling}
\end{figure*}
